% answers_short.tex —— 2010 年·普通高等学校招生考试·浙江卷·理科 · 简短答案速查
% 由 tools/make_short_answers.py 自 answers.tex 的 \gksolhead/\gkanswer 提取，
% 并与 2010/2010_数学真题_浙江理_答案要点.md 逐题比对后生成，勿手改（改源 MD 后重跑）。

\begin{gknotice}
\textbf{使用说明}：本节只给最终答案便于快速核对，完整推导见解析版。答案为编写者独立推导，\textbf{非官方标准答案}。
\end{gknotice}

\gkpart{选择题}
\begin{center}\begin{tabular}{ccccccccccc}
\toprule
题号 & 1 & 2 & 3 & 4 & 5 & 6 & 7 & 8 & 9 & 10 \\
\midrule
答案 & B & A & D & B & D & B & C & C & A & B \\
\bottomrule
\end{tabular}\end{center}

\gkpart{填空题}
\begin{itemize}[leftmargin=2.6em]
  \item[\textbf{11.}] $\pi$。
  \item[\textbf{12.}] $144$。
  \item[\textbf{13.}] $\dfrac{3\sqrt2}4$。
  \item[\textbf{14.}] $T_n=\begin{cases}0,&n\text{ \text{为偶数}},\\[2pt]\dfrac1{2^n}-\dfrac1{3^n},&n\text{ \text{为奇数}}.\end{cases}$
  \item[\textbf{15.}] $(-\infty,-2\sqrt2]\cup[2\sqrt2,+\infty)$。
  \item[\textbf{16.}] $\left(0,\dfrac{2\sqrt3}3\right]$。
  \item[\textbf{17.}] $264$。
\end{itemize}

\gkpart{解答题}
\begin{itemize}[leftmargin=2.6em]
  \item[\textbf{18.}] （1）$\dfrac{\sqrt{10}}4$；（2）$b=\sqrt6$，$c=4$ 或 $b=2\sqrt6$，$c=4$。
  \item[\textbf{19.}] （1）$P(\xi=0.5)=\dfrac14$，$P(\xi=0.7)=\dfrac14$，$P(\xi=0.9)=\dfrac12$，$E\xi=\dfrac34$；（2）$\dfrac38$。
  \item[\textbf{20.}] （1）$\dfrac{\sqrt3}3$；（2）$\dfrac{21}4$。
  \item[\textbf{21.}] （1）$x-\sqrt2y-1=0$；（2）$1<m<2$。
  \item[\textbf{22.}] （1）$b<-a$；（2）存在，共四组：$b=-a-3$ 时 $x_4=a+2\sqrt6$ 或 $a-2\sqrt6$；$b=-a-\dfrac{7+\sqrt{13}}2$ 时 $x_4=a+\dfrac{1+\sqrt{13}}2$；$b=-a-\dfrac{7-\sqrt{13}}2$ 时 $x_4=a+\dfrac{1-\sqrt{13}}2$。
\end{itemize}
